\clearpage
\fancyhead[L]{\small CARLINI. MAILÄNDER EPHEMERIDEN FÜR 1816.}
\SourceLeaf{597}
\newcommand{\ObservedLibelleDenominator}{\hbox{\vrule width0pt height4.600000bp depth0pt\special{pdf:bcontent}\special{pdf:code q 0 g 2.146667 4.600000 m 2.146667 4.293333 l 2.453333 3.986667 l 2.453333 3.680000 l 1.840000 3.066667 l 1.533333 3.066667 l 0.613333 2.146667 l 0.000000 2.146667 l 0.000000 0.613333 l 0.306667 0.613333 l 0.613333 0.306667 l 1.226667 0.306667 l 1.533333 0.000000 l 3.066667 0.000000 l 3.986667 0.920000 l 3.986667 1.840000 l 3.066667 2.760000 l 3.066667 3.066667 l 3.680000 3.680000 l 3.680000 3.986667 l 3.066667 4.600000 l h 1.533333 2.146667 m 1.840000 2.453333 l 2.453333 1.840000 l 2.453333 1.533333 l 1.840000 0.920000 l 1.533333 1.226667 l h f* Q}\special{pdf:econtent}\kern3.986667bp}}

\begin{center}\small Göttingische gelehrte Anzeigen. Stück 47. Seiten 457..46c. 23. März 1816.\\[3pt]\rule{.72\linewidth}{.3pt}\end{center}

\emph{Effemeridi astronomiche di Milano per l’anno bisestile} 1816, \emph{calcolate da} \textsc{Francesco Carlini}. \emph{Con appendice. Dalla cesarea regia stamperia.} Mailand 1815. Oktavformat. Die Ephemeriden umfassen 108 Seiten, der Anhang 100.

Die Beigaben zum astronomischen Kalender sind diesmal 1) das Verzeichnis der 34 Fundamentalsterne mit ihren jährlichen Präzessionen und den Konstanten zur Berechnung ihrer Aberration und Nutation, wie im Jahrgang 1815 (siehe unsere Anzeigen vom Vorjahr, S. 89 [S. 574 dieses Bandes]). 2) Die Refraktionstafel von \textsc{Carlini}. Diese Tafel, die auf der Hypothese von \textsc{Laplace} über den Zustand der Atmosphäre und der daraus abgeleiteten Theorie der Refraktion beruht, bei der jedoch die dabei auftretenden Konstanten anhand eigener Beobachtungen neu bestimmt wurden, war bereits im Jahrgang 1808 der Ephemeriden abgedruckt. Da der Rezensent diesen Jahrgang nicht zur Hand hat, kann er nicht angeben, ob die vorliegende Tafel bloß ein unveränderter Abdruck ist. 3) Verzeichnis der Bedeckungen von Fixsternen im Jahr 1816, berechnet von den Florentiner Astronomen für den Meridian und die Polhöhe von Florenz.

Der Anhang enthält folgende Aufsätze. 1) Beobachtete Refraktion bei kleinen Höhen über dem Horizont, von \textsc{Barnabas Oriani}. Unter den in den vorhergehenden Jahrgängen bekannt gemachten Beobachtungen mit dem Repetitionskreis von \textsc{Reichenbach} mit drei Fuß Durchmesser finden sich 19 im Jahr 1811 beobachtete untere Kulminationen des Sterns Capella, der in Mailand in einer Höhe von $1^\circ\;36'$, nur wenige Bogenminuten über den Bergen, die den Horizont begrenzen, durchgeht. Überdies waren im Jahr 1811 noch 14 obere Kulminationen beobachtet worden, aus denen mit der Polhöhe $45^\circ\;28'\;0''7$ die mittlere Deklination $45^\circ\;47'\;28''5$ für den Anfang des Jahres 1811 abgeleitet wurde. Die unteren Kulminationen wurden von Herrn \textsc{Oriani} nach den wichtigsten Refraktionstafeln berechnet, nämlich den Tafeln von \textsc{Bradley}, \textsc{Tob. Mayer}, \textsc{Piazzi}, \textsc{Delambre}, \textsc{Carlini} und \textsc{Bessel}, unter denen die von \textsc{Carlini} am besten und bis auf wenige Bogensekunden übereinstimmten, sofern nämlich die vermeintliche Verbesserung für die nördliche Hälfte des Meridians, die \textsc{Carlini} früher für nötig gehalten hatte, weggelassen wurde. 2) Fortsetzung der Beobachtungen über das periodische Schwanken der Gebäude, von \textsc{Angelo Cesaris}. Die schon im Jahrgang der Ephemeriden für 1813 mitgeteilten Erfahrungen über diesen Gegenstand (siehe unsere Anzeigen 1813, S. 951 [S. 556 d. B.]) werden hier um andere, nicht weniger bemerkenswerte Erfahrungen ergänzt. Dort war von einer Bewegung die Rede, die sich gleichsam als Rotation der Mauern, an denen Mauerquadrant und Mittagsfernrohr aufgehängt sind, um eine vertikale Achse auffassen lässt. Die neuen Erfahrungen beziehen sich auf eine Rotation um eine horizontale, von Osten nach Westen gerichtete Achse. Diese Bewegung wurde durch eine ausgezeichnete Libelle von \textsc{Reichenbach} sichtbar gemacht, die am Mauerquadranten angebracht und so empfindlich war, dass eine Änderung um eine Bogensekunde einen Ausschlag von $1\frac{2}{\ObservedLibelleDenominator}$ Linien ergab. Dieselben Veränderungen, die sich an dieser Libelle zeigten, ließen sich auch an der mit dem Mauerquadranten beobachteten Zenitdistanz sowie an einer zweiten Libelle erkennen, die an der gegenüberliegenden Seite der Mauer angebracht war. Die hier mitgeteilten Erfahrungen deuten sehr bestimmt auf einen Zusammenhang mit dem Wetter und mit der Stärke der Erwärmung des Gebäudes durch das Sonnenlicht. Am regelmäßigsten sind die Veränderungen vom Vormittag bis zum Nachmittag. Bei heiterem Wetter bewegt sich die Blase in der Libelle nach Süden. Die Bewegung beträgt im Durchschnitt etwa $2''$ und bleibt bei bedecktem Himmel ganz aus. 3) Tafel für die Mittelpunktsgleichung der Vesta (Exzentrizität $=0.0889$) und \SourceLeaf{598}für die Reduktion auf die Ekliptik (Neigung der Bahn $=7^\circ\;8'\;20''$, von \textsc{Carlini}. 4) Sternbedeckungen, beobachtet in Mailand in den Jahren 1811--1815 von ihm und in Florenz in den Jahren 1810--1813 von den dortigen Astronomen. --- 5) Schiefe der Ekliptik aus den Beobachtungen der Sonnenwenden mit dem Kreis von \textsc{Reichenbach} mit drei Fuß Durchmesser, von \textsc{Barnabas Oriani}. Die Beobachtungen bei der Wintersonnenwende 1810 und den beiden Sonnenwenden 1811 wurden bereits in früheren Jahrgängen mitgeteilt. Hier erhalten wir die ausführlichen Beobachtungen bei den fünf folgenden Sonnenwenden und die Ergebnisse, die aus allen Beobachtungen gewonnen wurden. Es ist sehr interessant, nun auch zu erfahren, welche Antwort die größeren Repetitionskreise von \textsc{Reichenbach} auf diese viel diskutierte Frage geben. Die vier Wintersonnenwenden ergeben für die auf den Anfang des Jahres 1812 reduzierte mittlere Schiefe der Ekliptik $=23^\circ\;27'\;48''20$, die vier Sommersonnenwenden $23^\circ\;27'\;50''77$. Es bleibt also ein Unterschied, der allerdings viel kleiner ist als derjenige, den die zwölfzölligen Kreise von \textsc{Reichenbach} und andere Instrumente ergeben haben. Zur Reduktion der Beobachtungen wurden die Refraktionstafeln von \textsc{Carlini} verwendet. Mit den Refraktionstafeln von \textsc{Delambre} wäre der Unterschied etwas größer und betrüge $3''71$. --- 6) Meteorologische Beobachtungen im Jahr 1814 von \textsc{Angelo Cesaris}, mit einer Übersicht über die Regenmengen in den Jahren 1764--1814. Aus allen 51 Jahren ergibt sich eine mittlere Regenmenge von 35 Zoll 3,92 Linien. Die Mengen der einzelnen Jahre schwanken natürlich sehr unregelmäßig. Fasst man sie aber für eine längere Reihe von Jahren zusammen, zeigt sich insgesamt eine stetige Zunahme. \textsc{Cesaris} schreibt dies der immer weiter verbreiteten Bewässerung der Felder zu. Dadurch werde eine große Menge Wasser über eine viel größere Fläche verteilt und stärkerer Verdunstung ausgesetzt, sodass sie in Form häufigeren Regens zurückkehren müsse. Die Zahl der Jahre ist jedoch wohl noch zu klein, und das Ergebnis hängt zu sehr von zufälligen Schwankungen ab, als dass sich \mbox{sichere} allgemeine Schlüsse ziehen ließen. Besonders feucht war das letzte Jahr 1814, in dem 58 Zoll 11,58 Linien Regen fielen. Die geringste Menge betrug im Jahr 1771 25 Zoll 11,6 Linien.
